\documentclass[11pt]{article}
\usepackage[margin=1in]{geometry}
\usepackage{xcolor}
\usepackage{tcolorbox}
\usepackage{hyperref}
\usepackage{enumitem}
\setlist{noitemsep,topsep=0pt}
\usepackage{booktabs}
\usepackage{array}
\usepackage{listings}
\usepackage{amsmath}

% Define OSU orange and AI blue colors
\definecolor{osaorange}{RGB}{220,115,38}
\definecolor{aiblue}{RGB}{30,144,255}
\definecolor{answergreen}{RGB}{0,128,0}

% Configure hyperref colors
\hypersetup{
    colorlinks=true,
    linkcolor=blue,
    urlcolor=blue,
    citecolor=blue
}

% Configure R code listings
\lstset{
    language=R,
    basicstyle=\ttfamily\small,
    breaklines=true,
    breakatwhitespace=true,
    postbreak=\mbox{\textcolor{red}{$\hookrightarrow$}\space},
    frame=single,
    backgroundcolor=\color{gray!10},
    numbers=left,
    numberstyle=\tiny,
    keywordstyle=\color{blue},
    commentstyle=\color{green!60!black},
    stringstyle=\color{red},
    columns=flexible,
    keepspaces=true,
    xleftmargin=0pt,
    xrightmargin=0pt,
    linewidth=\textwidth,
    aboveskip=\baselineskip,
    belowskip=\baselineskip,
    framexleftmargin=0pt,
    framexrightmargin=0pt
}

\title{}
\author{}
\date{}

\begin{document}

% Orange header box with black outline
\begin{tcolorbox}[colback=osaorange,colframe=black,boxrule=2pt,arc=0pt,left=6pt,right=6pt,top=10pt,bottom=10pt]
\centering
\textcolor{white}{\textbf{\Large In-Class Worksheet - ANSWER KEY}}\\
\textcolor{white}{\textbf{\large Introduction to R and Data Presentation}}
\end{tcolorbox}

\vspace{8pt}

\noindent\textbf{Course:} H524 Introduction to Biostatistics\\
\textbf{Topic:} Introduction to R, Data Presentation, Numerical Summary Measures

\vspace{8pt}

\section{Part 1: Getting Started with R}

\subsection{Exercise 1.1: Basic R Operations}
Work with a partner to complete these basic R calculations:

\begin{lstlisting}
# 1. Use R as a calculator
3 * 2 + 4
# Answer: 10

# 2. Create a simple vector of systolic blood pressure readings
blood_pressure <- c(170, 193, 110, 135, 135, 115)
# This creates a numeric vector with 6 blood pressure values
\end{lstlisting}

\textbf{Discussion Questions - Answers:}
\begin{itemize}
\item \textcolor{answergreen}{\textbf{How do you create a vector in R?}} Use the \texttt{c()} function with values separated by commas
\item \textcolor{answergreen}{\textbf{What does the \texttt{c()} function do?}} It combines (concatenates) values into a vector
\end{itemize}

\section{Part 2: Data Types and Classification}

\subsection{Exercise 2.1: Identifying Data Types}
For each variable below, classify it as:
\begin{itemize}
\item \textbf{Nominal} (categorical without order)
\item \textbf{Ordinal} (categorical with order)
\item \textbf{Discrete} (countable numbers)
\item \textbf{Continuous} (measurable numbers)
\end{itemize}

\begin{center}
\begin{tabular}{|>{\raggedright\arraybackslash}p{4.8cm}|>{\centering\arraybackslash}p{2.2cm}|>{\raggedright\arraybackslash}p{5.5cm}|}
\hline
\textbf{Variable} & \textbf{Type} & \textbf{Reasoning} \\
\hline
Blood type (A, B, AB, O) & \textcolor{answergreen}{\textbf{Nominal}} & \textcolor{answergreen}{Categories with no natural order} \\[2ex]
\hline
Pain level (none, mild, moderate, severe) & \textcolor{answergreen}{\textbf{Ordinal}} & \textcolor{answergreen}{Categories with clear ordering from least to most severe} \\[2ex]
\hline
Number of hospitalizations this year & \textcolor{answergreen}{\textbf{Discrete}} & \textcolor{answergreen}{Countable whole numbers (0, 1, 2, 3, ...)} \\[2ex]
\hline
Body temperature (°F) & \textcolor{answergreen}{\textbf{Continuous}} & \textcolor{answergreen}{Can be measured to any level of precision} \\[2ex]
\hline
Treatment group (control, low dose, high dose) & \textcolor{answergreen}{\textbf{Nominal}} & \textcolor{answergreen}{Categories with no inherent order} \\[2ex]
\hline
\end{tabular}
\end{center}

\subsection{Exercise 2.2: AI Enhancement Activity}
\begin{tcolorbox}[colback=aiblue!10,colframe=aiblue,boxrule=1pt,arc=2pt,left=6pt,right=6pt,top=6pt,bottom=6pt]
\textbf{\textcolor{aiblue}{Using AI Tools:}} Use GitHub Copilot or another AI assistant to help you:

\begin{enumerate}
\item Generate R code to create vectors for each data type above
\item Ask AI: ``What's the best way to store categorical data in R?''
\end{enumerate}

\textbf{Sample AI-Generated Code:}
\begin{lstlisting}
# Nominal data - blood type
blood_type <- factor(c("A", "B", "AB", "O"))

# Ordinal data - pain level
pain_level <- factor(c("none", "mild", "moderate", "severe"),
                     levels = c("none", "mild", "moderate", "severe"),
                     ordered = TRUE)

# Discrete data - hospitalizations
hospitalizations <- c(0, 1, 2, 1, 0, 3)

# Continuous data - temperature
temperature <- c(98.6, 99.2, 97.8, 100.1, 98.9)

# Nominal data - treatment group
treatment <- factor(c("control", "low dose", "high dose"))
\end{lstlisting}

\textbf{AI Answer:} Use \texttt{factor()} for categorical data, with \texttt{ordered = TRUE} for ordinal data.
\end{tcolorbox}

\newpage
\section{Part 3: Descriptive Statistics}

\subsection{Exercise 3.1: Measures of Central Tendency}
Using the blood pressure data: \texttt{blood\_pressure <- c(170, 193, 110, 135, 135, 115)}

Complete the calculations both by hand and in R:

\begin{lstlisting}
# Create the data
blood_pressure <- c(170, 193, 110, 135, 135, 115)

# Calculate mean
mean_bp <- mean(blood_pressure)
mean_bp  # 143

# Calculate median
median_bp <- median(blood_pressure)
median_bp  # 135

# Find the mode (most frequent value)
table(blood_pressure)  # Shows 135 appears twice
# Mode: 135

# Order the data
ordered_bp <- sort(blood_pressure)
ordered_bp  # 110 115 135 135 170 193
\end{lstlisting}

\textbf{Hand Calculations:}
\begin{itemize}
\item Mean: \textcolor{answergreen}{\textbf{143}} mmHg \textcolor{answergreen}{(170+193+110+135+135+115)/6 = 858/6 = 143}
\item Median: \textcolor{answergreen}{\textbf{135}} mmHg \textcolor{answergreen}{(135+135)/2 = 135 (middle values)}
\item Mode: \textcolor{answergreen}{\textbf{135}} mmHg \textcolor{answergreen}{(appears twice)}
\end{itemize}

\textbf{R Results:}
\begin{itemize}
\item Mean: \textcolor{answergreen}{\textbf{143}} mmHg
\item Median: \textcolor{answergreen}{\textbf{135}} mmHg
\item Mode: \textcolor{answergreen}{\textbf{135}} mmHg
\end{itemize}

\newpage
\section{Part 4: Measures of Variability}

\subsection{Exercise 4.1: Spread of Data}
\begin{lstlisting}
# Create the data
blood_pressure <- c(170, 193, 110, 135, 135, 115)

# Calculate range
max_bp <- max(blood_pressure)  # 193
min_bp <- min(blood_pressure)  # 110
range_bp <- max_bp - min_bp    # 83

# Calculate variance
var_bp <- var(blood_pressure)  # 1116.4

# Calculate standard deviation
sd_bp <- sd(blood_pressure)    # 33.41

# Calculate standard error of the mean
n <- length(blood_pressure)    # 6
sem_bp <- sd_bp / sqrt(n)      # 13.64
\end{lstlisting}

\textbf{Fill in your results:}
\begin{itemize}
\item Range: \textcolor{answergreen}{\textbf{83}} mmHg \textcolor{answergreen}{(193 - 110)}
\item Variance: \textcolor{answergreen}{\textbf{1116.4}} mmHg$^2$ \textcolor{answergreen}{(sample variance)}
\item Standard Deviation: \textcolor{answergreen}{\textbf{33.41}} mmHg \textcolor{answergreen}{(square root of variance)}
\item Standard Error of Mean: \textcolor{answergreen}{\textbf{13.64}} mmHg \textcolor{answergreen}{(33.41/$\sqrt{6}$)}
\end{itemize}

\textcolor{answergreen}{\textbf{Discussion:}} These measures tell us the blood pressure data has considerable spread (range = 83), high variability (SD = 33.41), and the mean has substantial uncertainty (SEM = 13.64).

\newpage
\section{Part 5: AI-Enhanced Data Exploration}

\subsection{Exercise 5.1: Using AI for Data Summary}
\begin{tcolorbox}[colback=aiblue!10,colframe=aiblue,boxrule=1pt,arc=2pt,left=6pt,right=6pt,top=6pt,bottom=6pt]
\textbf{\textcolor{aiblue}{Prompt for AI:}} ``Help me create a comprehensive summary of this blood pressure dataset: c(170, 193, 110, 135, 135, 115). Include all relevant descriptive statistics and suggest appropriate visualizations.''
\end{tcolorbox}

\textbf{Sample AI-Generated Code:}
\begin{lstlisting}
# Blood pressure data analysis
blood_pressure <- c(170, 193, 110, 135, 135, 115)

# Comprehensive summary
summary(blood_pressure)
#   Min. 1st Qu.  Median    Mean 3rd Qu.    Max.
#    110     125     135     143     158     193

# Additional statistics
length(blood_pressure)    # n = 6
sd(blood_pressure)        # SD = 33.41
var(blood_pressure)       # Var = 1116.4
range(blood_pressure)     # Range: 110 193
IQR(blood_pressure)       # IQR = 33

# Suggested visualizations
hist(blood_pressure, main="Blood Pressure Distribution",
     xlab="Systolic BP (mmHg)", col="lightblue")

boxplot(blood_pressure, main="Blood Pressure Boxplot",
        ylab="Systolic BP (mmHg)")

dotchart(blood_pressure, main="Blood Pressure Dot Plot")
\end{lstlisting}

\textbf{Verification Steps - Sample Answers:}
\begin{itemize}
\item \textcolor{answergreen}{\textbf{Run the code and check if results match:}} Yes, all calculations match hand calculations
\item \textcolor{answergreen}{\textbf{Do the suggested visualizations make sense?}} Yes, histogram shows distribution, boxplot shows quartiles and outliers, dot plot shows individual values
\item \textcolor{answergreen}{\textbf{Are there any errors or improvements needed?}} Code is correct, might add axis labels for clarity
\end{itemize}

\textbf{What you learned from AI - Sample Answers:}
\begin{itemize}
\item New R functions discovered: \textcolor{answergreen}{\textbf{summary(), IQR(), dotchart()}}
\item Visualization suggestions: \textcolor{answergreen}{\textbf{histogram, boxplot, dot chart}}
\item Any surprising insights: \textcolor{answergreen}{\textbf{IQR = 33, shows middle 50\% spread}}
\end{itemize}

\section{Wrap-up Discussion}

\textbf{Reflection Questions - Sample Answers:}
\begin{enumerate}
\item \textcolor{answergreen}{\textbf{When might you prefer mean vs. median?}} Use median when data is skewed or has outliers; use mean when data is normally distributed
\item \textcolor{answergreen}{\textbf{How can AI tools help while ensuring accuracy?}} AI generates code quickly, but always verify results by hand or with known datasets; check for logical errors
\item \textcolor{answergreen}{\textbf{What's one new R function you learned?}} Examples: \texttt{factor()}, \texttt{summary()}, \texttt{IQR()}, \texttt{table()}
\end{enumerate}

\textbf{Next Class Preview:}
We'll explore data visualization with histograms and boxplots, plus dive deeper into R programming basics.

\section{AI Usage Documentation - Sample}
\begin{tcolorbox}[colback=gray!10,colframe=gray,boxrule=1pt,arc=2pt,left=6pt,right=6pt,top=6pt,bottom=6pt]
\textbf{Sample documentation:}
\begin{itemize}
\item AI tools used: \textcolor{answergreen}{\textbf{GitHub Copilot, ChatGPT}}
\item Key prompts: \textcolor{answergreen}{\textbf{"Generate R code for data summary", "Best way to store categorical data"}}
\item Verification methods: \textcolor{answergreen}{\textbf{Hand calculations, checking output against expected results}}
\item Errors found and corrected: \textcolor{answergreen}{\textbf{None in this exercise, but always double-check factor levels}}
\end{itemize}
\end{tcolorbox}

\end{document}
